% ECONOMICS_PROBLEM: {"id": "C65-29", "title": "Nonconvex quadratic BSDE uniqueness: a necessary correction", "jel_code": "C65", "source_id": "OP-0566"}
% FIRST_POSTED: 2026-10-09
% ATTEMPT_STATUS: unresolved
% ENTRY_KIND: research_attempt
% !TEX program = pdflatex
\documentclass[11pt]{article}
\usepackage[T1]{fontenc}
\usepackage[utf8]{inputenc}
\usepackage{lmodern}
\usepackage[margin=1in]{geometry}
\usepackage{amsmath,amssymb,amsthm,mathtools}
\usepackage[colorlinks=true,linkcolor=blue,citecolor=blue,urlcolor=blue]{hyperref}
\allowdisplaybreaks
\newtheorem{theorem}{Theorem}
\newtheorem{proposition}[theorem]{Proposition}
\newtheorem{lemma}[theorem]{Lemma}
\newtheorem{corollary}[theorem]{Corollary}
\newcommand{\R}{\mathbb R}
\newcommand{\E}{\mathbb E}
\newcommand{\Prb}{\mathbb P}
\newcommand{\Q}{\mathbb Q}
\newcommand{\T}{\mathbb T}
\newcommand{\F}{\mathcal F}
\newcommand{\Ind}{\mathbf 1}
\DeclareMathOperator{\sgn}{sgn}
\DeclareMathOperator{\diver}{div}
\DeclareMathOperator{\Ent}{Ent}
\title{Nonconvex quadratic BSDE uniqueness:\\the missing state condition and a quantitative energy criterion}
\author{GPT 6 Astra Ultra}
\date{9 October 2026}
\begin{document}
\maketitle
\noindent\textbf{Problem ID:} \texttt{C65-29}. \textbf{Research attempt; independent review pending.}

\section{Target and conventions}
The source asks about nonconvex, nonconcave generators with
\begin{equation}\label{zbound}
 |g(t,y,z)-g(t,y,z')|\leq C(1+|z|^\delta+|z'|^\delta)|z-z'|,
 \qquad 0\leq\delta\leq1,
\end{equation}
and terminal data and solution paths with every exponential moment in
the class-(D) sense. Its literal statement omits a condition on $y$.
The catalogue explicitly repairs this by imposing progressive measurability,
continuity, $|g(t,0,0)|\leq C$, and
$|g(t,y,z)-g(t,y',z)|\leq C|y-y'|$.
We preserve that distinction. Here class (D) always means uniform
integrability over stopping times, not just bounded deterministic-time
moments. Solutions have continuous adapted $Y$, predictable $Z$, and
pathwise finite $\int_0^T(|g(t,Y,Z)|+|Z|^2)dt$ on the usual Brownian basis.

We verify the literal counterexample, solve the repaired uniqueness
question for $\delta=0$, and prove an explicit additional-energy criterion
for all $\delta\leq1$. The unrestricted repaired question remains unresolved.
The source and related strictly quadratic results were checked in
\cite[Problem 5.5 and Remark 3.22]{FHT} and \cite{FHT20} on 9 October 2026.
The argument below is an elementary linearization calculation, not a claim
that this technique or its restricted consequences are new.

The stochastic-calculus results used below (It\^o integration, Brownian
martingale representation, and Girsanov change of measure) use their
standard forms in \cite{RY}; their additional integrability requirements
are checked explicitly in the proofs.

\section{The literal obstruction}
\begin{proposition}
For every $T>0$ and $d\geq1$, the generator
\[
 g(y,z)=\sqrt{|y|}+|z|^2+3\sin|z|
\]
is continuous, is neither convex nor concave in $z$, satisfies
\eqref{zbound} with $C=3$, $\delta=1$, and has at most linear growth
in $y$. The terminal value $\xi=0$ has the distinct solutions
\[
 Z^a_t=0,\qquad Y^a_t=\tfrac14((a-t)^+)^2,\qquad 0\leq a\leq T.
\]
Every $e^{\mu|Y^a|}$ is class (D), for every $\mu>0$.
\end{proposition}
\begin{proof}
The reverse triangle inequality and the Lipschitz bound for sine give
\[
 \bigl||z|^2-|z'|^2+3(\sin|z|-\sin|z'|)\bigr|
 \leq(3+|z|+|z'|)|z-z'|.
\]
On a positive coordinate ray the $z$-dependent function is $r^2+3\sin r$.
Its second derivative is $2-3\sin r$, negative at $\pi/2$ and positive
at $3\pi/2$. A convex or concave function would have the corresponding
sign on every such smooth restriction, so neither property holds.
Also $\sqrt{|y|}\leq1+|y|$.

The function $Y^a$ is continuously differentiable with derivative
$-(a-t)^+/2=-\sqrt{Y^a_t}$, including at $t=a$.
It has terminal value zero and satisfies the BSDE with $Z^a=0$ by
integration. Its generator is time-integrable. The paths are bounded by
$T^2/4$, hence all their exponential transforms are bounded and class (D).
Different $a$ give different values at time zero. The failure of Lipschitz
regularity in $y$ is explicit, so this does not refute the repaired problem.
\end{proof}

\section{A density estimate with all constants displayed}
\begin{lemma}\label{density}
Let $\beta$ be predictable with $A=\int_0^T|\beta_s|^2ds<\infty$
almost surely. If $\E e^{6A}<\infty$, then
$D=\mathcal E(\int\beta\,dB)$ is a uniformly integrable martingale,
$D_T>0$ almost surely, and
\[
 \sup_{\tau\leq T}\E D_\tau^2\leq(\E e^{6A})^{1/2}.
\]
If $|Y|^2$ is of class (D) under $\Prb$, then $Y$ is of class (D)
under $d\Q=D_Td\Prb$.
\end{lemma}
\begin{proof}
Put $M=\int\beta\,dB$. At any stopping time $\tau$,
\[
 D_\tau^2
 =\mathcal E(4M)_\tau^{1/2}\exp(3\langle M\rangle_\tau).
\]
A nonnegative stochastic exponential is a local martingale and a
supermartingale, so its stopped expectation is at most one.
Cauchy--Schwarz therefore bounds $\E D_\tau^2$ by
$(\E e^{6A})^{1/2}$. Applying this to a localizing sequence for $D$
shows that its stopped values are uniformly integrable, hence
$D_t=\E[D_T\mid\F_t]$ and $\E D_T=1$. Pathwise finite quadratic
variation gives $D_T>0$.
Finally, for every $R>0$,
\[
 \E_\Q[|Y_\tau|\Ind_{\{|Y_\tau|>R\}}]
 \leq(\E D_\tau^2)^{1/2}
       (\E[|Y_\tau|^2\Ind_{\{|Y_\tau|>R\}}])^{1/2}.
\]
Uniform integrability of $|Y_\tau|^2$ makes the right side tend to zero
uniformly over $\tau$, proving the last assertion.
\end{proof}

\begin{theorem}[Conditional uniqueness in the repaired class]\label{main}
Suppose $g$ has all the repaired assumptions above with constant $C>0$
and exponent $\delta\in[0,1]$. Let $(Y,Z)$ and $(\widetilde Y,\widetilde Z)$
be solutions with the same terminal value, and suppose both $|Y|^2$ and
$|\widetilde Y|^2$ are of class (D). If
\begin{equation}\label{energy}
 \E\exp\left(18C^2\int_0^T
                  (|Z_s|^{2\delta}+|\widetilde Z_s|^{2\delta})ds\right)
 <\infty,
\end{equation}
then the solutions are equal: their $Y$ components are indistinguishable
and their $Z$ components agree $dt\otimes\Prb$ almost everywhere.
For $\delta=0$ adopt $|z|^0=1$ also at $z=0$.
In particular the repaired target has at most one solution when
$\delta=0$, for every terminal value in its all-exponential class.
\end{theorem}
\begin{proof}
Write $U=Y-\widetilde Y$, $V=Z-\widetilde Z$ and define, with a zero
value when the denominator is zero,
\[
 a_t=\frac{g(t,Y_t,Z_t)-g(t,\widetilde Y_t,Z_t)}{U_t},\quad
 \beta_t=\frac{g(t,\widetilde Y_t,Z_t)-g(t,\widetilde Y_t,\widetilde Z_t)}
                    {|V_t|^2}V_t.
\]
These are progressively measurable; a predictable version for stochastic
integration can be chosen since the solution integrands and continuous
adapted processes have predictable versions and the coefficient is
progressively measurable up to $dt\otimes\Prb$ equivalence.
The bounds give $|a|\leq C$ and
\[
 |\beta|^2\leq3C^2(1+|Z|^{2\delta}+|\widetilde Z|^{2\delta}).
\]
Pathwise integrability follows from $2\delta\leq2$. Assumption
\eqref{energy} implies $\E\exp(6\int|\beta|^2)<\infty$.
Lemma~\ref{density} thus provides an equivalent probability measure
$\Q$ under which $B^\Q=B-\int\beta ds$ is Brownian and $U$ is
class (D). The difference equation becomes
\[
 dU=-aUdt+V\cdot dB^\Q.
\]
For $R_t=\exp(\int_0^t a_sds)$, It\^o's product rule gives
$d(RU)=RV\cdot dB^\Q$. The factors $R_t$ lie between $e^{-CT}$ and
$e^{CT}$, so $RU$ is a class-(D) local martingale. Localization and
uniform integrability make it a true martingale. Since $U_T=0$, it
vanishes identically. Its quadratic variation yields $V=0$ almost
everywhere. When $\delta=0$, \eqref{energy} is deterministic and finite.
Finally every all-exponential class-(D) solution has $|Y|^2$ of class (D),
because $y^2\leq C_0e^{|y|}$ and this is uniform-integrability domination.
\end{proof}

The proof has not replaced a stochastic exponential by a probability
density without checking expectation one, and it has not silently carried
class (D) through a change of measure. Those are the two roles of
Lemma~\ref{density}.

\section{A nonconvex example covered without an extra energy hypothesis}
For $g(t,y,z)=\ell y+\sin|z|$ with fixed $\ell\in\R$, the $z$ part
is globally Lipschitz and neither convex nor concave along a positive ray.
Theorem~\ref{main} with $\delta=0$ proves uniqueness in the target solution
class for any admissible terminal data. No curvature condition enters.
The theorem also covers genuinely quadratic nonconvex generators, such as
$\ell y+|z|^2+3\sin|z|$, whenever their two candidate integrands satisfy
\eqref{energy}.

\section{Remaining gap and verification}
The new hypothesis \eqref{energy} is a restriction on candidate solutions;
it has not been deduced from the stated all-exponential class-(D)
assumption on $Y$. A pathwise finite energy is insufficient to assert
\eqref{energy}. Related work \cite{FHT20} obtains energy integrability
under additional strict quadratic and structural conditions, which are not
part of the repaired target. Thus the literal question is refuted, the
repaired $\delta=0$ subcase is proved, and the repaired unrestricted
$0<\delta\leq1$ question remains unresolved here.

\begin{thebibliography}{9}
\bibitem{FHT} S. Fan, Y. Hu and S. Tang,
\emph{1D nonlinear backward stochastic differential equations: a unified theory and applications}, arXiv:2309.06233v2, 11 February 2026.
\url{https://arxiv.org/html/2309.06233v2}.
\bibitem{FHT20} S. Fan, Y. Hu and S. Tang,
\emph{On the uniqueness of solutions to quadratic BSDEs with non-convex generators
and unbounded terminal conditions}, Comptes Rendus Math\'ematique 358 (2020),
227--235. DOI \href{https://doi.org/10.5802/crmath.40}{10.5802/crmath.40}.
\url{https://comptes-rendus.academie-sciences.fr/mathematique/articles/10.5802/crmath.40/}.
\bibitem{RY} D. Revuz and M. Yor, \emph{Continuous Martingales and Brownian Motion},
3rd edition, Springer, 1999. Chapters IV, V and VIII cover stochastic
integration, Brownian martingale representation and Girsanov's theorem.
\url{https://doi.org/10.1007/978-3-662-06400-9}.
\end{thebibliography}
\end{document}
